% Created 2019-05-31 Fri 17:25
% Intended LaTeX compiler: pdflatex
\documentclass[11pt]{teraum_adventure}
\usepackage[utf8]{inputenc}
\usepackage[T1]{fontenc}
\usepackage{graphicx}
\usepackage{grffile}
\usepackage{longtable}
\usepackage{wrapfig}
\usepackage{rotating}
\usepackage[normalem]{ulem}
\usepackage{amsmath}
\usepackage{textcomp}
\usepackage{amssymb}
\usepackage{capt-of}
\usepackage{hyperref}
\author{emsenn}
\date{\today}
\title{The Crown of Idbol Bark}
\hypersetup{
 pdfauthor={emsenn},
 pdftitle={The Crown of Idbol Bark},
 pdfkeywords={},
 pdfsubject={},
 pdfcreator={Emacs 26.2 (Org mode 9.2.3)},
 pdflang={English}}
\begin{document}

\section*{The Crown of Idbol Bark}
\label{sec:org01ad761}
\textbf{\emph{The Crown of Idbol Bark}} is a \emph{tabletop role-playing \textbf{adventure}} in which the party attempts to arrest an upstart "king", usually taking two to four hours to complete.
This adventure was written to be played with the included rules, but can be adapted for any system.

This adventure is ready to be used, but isn't finalized. Pay special attention to text flagged "TK." Things might be unfair to players, or unclear to narrators. If you encounter any issues, please send me an email at emsenn@emsenn.net

\noindent\rule{.3\textwidth}{0.8pt}

Danger is never far in the \emph{Gloaming}, where marshes and mesas give
home to bandits and beasts. The lone \emph{trade road} through the region
is a gauntlet for traveling merchants.

The town of \emph{Belcaer} is a haven for those travelers. The same
\emph{magic}-users who built the \emph{trade road} built the town as a fortress.
On top of a limestone mesa, the town now exists to serve travelers.
Residents offer meager services, hoping to win the travelers' coin.

Recently there haven't been many coins to take. A bandit calling
himself "King Brag" has been mugging travelers. Now, farmers are
reporting missing livestock.

A local official named \emph{Abeth Harbrook} believes she has found where
the "king" is hiding: \emph{Jacob's Folly}, an tower built and abandoned in
the forgotten past.

She's offering to pay anyone who captures the "king," and the party
has gathered at her office.

\textbf{\emph{}}

This adventure is set in tragically funny fantasy world of \href{../index.org}{Teraum}, in the year 80pc, during the \emph{Reconstruction Era}.
TK

This adventure comes with suggested characters that the players could use as their adventurers. See the "Character Templates" supplemental section.

To the extent possible under U.S. law, the author of this work has
waived all copyright and related or neighboring rights to this work.


There are other adventures set in Teraum, along with lots of other material! To view the entire library, visit \href{https://emsenn.net/teraum/}{https://emsenn.net/teraum/}

\noindent\rule{.3\textwidth}{0.8pt}

Before beginning the session, your party should know:
\begin{itemize}
\item How much does each of them weigh, with and without equipment?
\item Is anyone a high-ranking officer in the \emph{Chandlers Guild}?
\end{itemize}

\begin{center}\includegraphics[width=.4\linewidth]{../assets/teraum-logotype-bw.png}\end{center}
\begin{center}\includegraphics[width=.3\linewidth]{../assets/public-domain.png}\end{center}

\newpage

\textbf{\emph{NOTE:}} Don't read past this point unless you're planning to \emph{narrate} this adventure. If you aren't sure what that means, read the included rules.
\subsection*{Exposition}
\label{sec:org61264a9}
\subsubsection*{Settings}
\label{sec:orgc2ee6c9}
\begin{itemize}
\item \textbf{\emph{Belcaer:}} A town of about four-hundred, built of stone atop a
mesa, overlooking the only major road in the \emph{Gloaming}.
\item \textbf{\emph{The Gloaming:}} A region of fens and marshes, broken up by mesas
and rock outcroppings. The region is unsettled and lawless.
\item \textbf{\emph{Jacob's Folly:}} A three-storey limestone tower built about eight
miles west of \emph{Belcaer}, far from any road or notable landmark.
\end{itemize}
\subsubsection*{Cast}
\label{sec:org7244b9a}
\begin{itemize}
\item \textbf{\emph{Abeth Harbrook:}} The local \emph{Chandlers Guild} representative.
\item \textbf{\emph{King Idbol Bark:}} A \emph{goblin} who lives with his clan in \emph{Jacob's
Folly}.
\end{itemize}
\subsection*{The Plot}
\label{sec:orgce0ca4e}
\subsubsection*{Scene One: Receiving the Quest}
\label{sec:org1801121}
\index{Belcaer!"Scene One" in /The Crown of Idbol Bark/}
\index{Harbrook, Abeth!"Scene One" in /The Crown of Idbol Bark/}
The members of the party are to meet, together or individually, with
\emph{Abeth Harbrook}, at her offices in the \emph{Guildhouse}. She is the
highest-ranked local official of the \emph{Chandlers Guild}, a commercial
organization based in the distant city of \emph{Ack}.

They rely on trade coming through the \emph{Gloaming}, so have recently
invested in the region and the town. Much of the \emph{trade road} has been
recently covered in gravel. In town, they built the \emph{Guildhouse}, a
massive pine structure that sticks out over the edge of the mesa, held
up by a lattice of timber.

\emph{Abeth Harbrook} has kept the party waiting, and now comes out to meet
them in the waiting room. She apologies for addressing them as a
group, if they aren't already associated.

Two seasons ago, \emph{Abeth Harbrook} first heard rumors of a "king," from
travelers that claimed they were mugged on the \emph{trade road}. At first
she dismissed the claims as exaggerations. She thought they wanted to
get out of paying guild tariffs.

Last season, livestock began to disappear from local farms. Barns and
silos were painted to say "\emph{King Bark} Was Here" and "\emph{Idbol Bark} Is
King."

Recent attempts to arrest the "king" failed, but revealed the location
of his hideout: \emph{Jacob's Folly}, a short tower built eight miles
outside of town.

\textbf{\emph{How does the party plan to capture} King Idbol Bark?}

If the party \textbf{risks} interrogating \emph{Abeth Harbrook}, she may include
that victims claimed the bandits were exceptionally tall and wore
masks. She withheld this information because she thought it was an
exaggeration.

An adventurer \textbf{knowledgable} about \emph{goblin} culture may recognize
\emph{"Bark"} as a \emph{goblin} surname.
\subsubsection*{Scene Two: Travel to the Tower}
\label{sec:orge07939b}
\index{Gloaming!"Scene Two" in /The Crown of Idbol Bark/}
\emph{Jacob's Folly} is approximately eight miles from \emph{Belcaer}. Six of
those miles are northwest along the \emph{trade road}, The final two
are through shallow marsh, broken by massive limestone outcroppings.

\textbf{\emph{How does the party get to Jacob's Folly?}}

If traveling at night or in inclement weather, the party \textbf{risks}
getting lost, though \textbf{knowledge} of navigation would help.

If the party travels quietly (such as at night or in inclement
weather,) as they approach the tower, they will remain undetected
until reaching "Scene Three."

Otherwise, they will be met by a group of "bandits" on the road, who
declare, in sqeaky rough voices, that declares they are in the
territory of \emph{King Bark}, and must surrender their money and supplies.

Each bandit is actually two \emph{goblins} in a long cloak, one standing on
the others shoulders and wearing a mask.

\textbf{\emph{How does the party handle these "bandits"?}}

If the party has horses or other livestock, the goblins will attempt
to steal these even if the party refuses to surrender. Otherwise, they
are more inclined to run than fight. If captured, they will happily
give the party directions to \emph{Jacob's Folly}, saying that \emph{King Bark}
would be happy for them to come bend the knee.
\subsubsection*{Scene Three: The Tower Exterior}
\label{sec:org899897e}
\index{Jacob's Folly!"Scene Three" in /The Crown of Idbol Bark/}
The limestone outcroppings in the marsh hide \emph{Jacob's Folly} until the
party is nearly there.

The limestone tower is three-storeys, and approximately twenty-five
feet tall and ten feet across. It is in a clearing about fifty yards
across. The tower's entrance has been sealed with the same limestone
bricks as the tower. There are a couple of windows near the top. Pine
planks, five on each side, are shoved into the mortar. They form an
unstable ladder up to the windows.

\emph{Goblins} are working around the tower, tearing grass out of the muck
and stuffing it into sacks. There is a burnt-out bonfire with horse
bones in it.

\textbf{\emph{How does the party get into Jacob's Folly?}}

When a member of the party moves into the clearin,g any goblins near
them will quietly move toward the tower, informing others on the way.

Some goblins will remain outside the tower, hiding behind rock
outcroppings in the field and throwing stones and debris at the party.
If the party \textbf{risks} engaging with any of those \emph{goblins}, goblins at
the top of the tower will have time to start burning straw on the roof
the tower, which they will drop on the party when they approach the
bottom of the tower.

If approaching in the evening, the party \textbf{risks} being noticed by an
owlbear mother that nests on the top of the tower, who will swoop down
and harass them, making the \emph{goblins} aware of them. The \emph{owlbear
mother} will attempt to carry any small (<120lb) adventurer into its
nest, where two hungry \emph{owlbear hatchlings} are waiting for them.


Placing more than 120lb on any single rung of the goblins' improvised
ladder is \textbf{risk}, with every 20lb making it \textbf{riskier}. Depending on
the result, a rung could loosen or break. The adventurer might gain a
fear of heights, or could fall, and even break a bone.

An adventurer who gets to the roof of \emph{Jacob's Folly} would be able
to take those \emph{hatchlings} and raise them, but unless the \emph{mother} is
killed, she will follow her hatchlings to the ends of the world and
try and get them back.
\subsubsection*{Scene Four: The Tower Interior}
\label{sec:orgb3bc24a}
\index{Jacob's Folly!"Scene Four" in /The Crown of Idbol Bark/}
The inside of the tower is dark, stinky, and beginning to fill with
acrid smoke.

When the party enters the room at the top of the tower, through the
window on the third storey, they'll find an empty room with some
sticks, dry piles of dung, buckets, and a large brass dish that's
clearly been used for fires, and might still have straw burning in it.
The goblins will retreat down underneath the trap door on the ground
floor.

There are two other "floors" inside the tower, though they are empty
except for more dirty straw and dung. There is a large, heavy wood
door in the floor at the ground level.


Once the party has entered the tower, all the goblins are underneath
this trap door, and have removed the ladder normally placed under it.
To get to the floor below is a seven-foot drop.

Smoke rises through the trapdoor, filling the tower with thicker
smoke that makes every risk riskier.

If the party is on the ground floor during the daytime, they will see
light entering the room through cracks in the filled-in doorway. A
\textbf{strong} character might \textbf{risk} knocking the mortar down, letting the
smoke out of the tower.

Opening the trapdoor exposes the party to further harassment from the
goblins, but provides an opportunity for dialog. The goblins will
call the party cowards, and say that if they come down there, \emph{King
Bark} is going to get them.

\textbf{\emph{How does the party get to} King Bark?}
\subsubsection*{Scene Five: The Tower Basement}
\label{sec:org906715b}
\index{Jacob's Folly!"Scene Five" in /The Crown of Idbol Bark/}
\index{Bark, Idbol!"Scene Five" in /The Crown of Idbol Bark/}
\index{Crown of Idbol Bark!"Scene Five" in /The Crown of Idbol Bark/}
The room is circular, a bit wider in diameter than the tower above,
and is full of farm tools and preserved foodstuff and grains.

Among the approximately two-dozen goblins in this room is one wearing
a "crown" made of twisted-together silverware: \emph{Idbol Bark}, leader of
this group of goblins.

\textbf{\emph{How does the party capture King Bark?}}

The goblins have no further means of retreat, so will fight until
unable. King Bark has no exceptional combat skills, and no goblin is
armed with anything more than farm and kitchen tools.
\subsection*{Epilogue}
\label{sec:orgfc8b35b}
\index{Harbrook, Abeth!"Epilogue" in /The Crown of Idbol Bark/}
\index{Crown of Idbol Bark!"Epilogue" in /The Crown of Idbol Bark/}
Once \emph{King Bark} and his goblins have been captured, disbanded, or
eliminated, the party should return to \emph{Belcaer} and tell \emph{Abeth
Harbrook} what happened, and potentially turn over custody of King
Bark and other goblins to her custody.

\begin{itemize}
\item How will \emph{Abeth Harbrook} react to claims or evidence of \emph{goblins}?
\item Who has possession of the \emph{Crown of Idbol Bark}? Each utensil in it
was magic in some way, though they're inactive now. To the right
person, or in the right future, it might be quite valuable.
\end{itemize}
\section*{Supplements}
\label{sec:orga07c8ef}
\subsection*{Character Templates}
\label{sec:org373f73d}
\begin{itemize}
\item A \textbf{\emph{high-ranking investigator from the} Chandlers Guild.} Among other qualities, they're \textbf{skilled} at negotiation and \textbf{knowledgeable} about bureaucracy. They're
here because the Guild is upset with \emph{Abeth Harbrook} for not solving
the problem already.
\begin{itemize}
\item How do they feel about Ms. Harbrook's performance?
\end{itemize}
\item A \textbf{\emph{local rancher}}. Among other qualities, they're \textbf{skilled} at navigating the region and \textbf{knowledgeable} about local flora and fauna. They're here
because \emph{King Bark}'s associates stole their favorite horse.
\begin{itemize}
\item What was their horse's name?
\end{itemize}
\item An out-of-work \textbf{\emph{academic}}.
They're here because they've heard of strange happenings in the region, and where there's mystery, there's money! Among other qualities, they're \textbf{skilled} at lying and \textbf{knowledgeable} about reading and writing.
\begin{itemize}
\item How can they exploit the mystery of \emph{King Brag's} identity?
\end{itemize}
\end{itemize}
\subsection*{Index}
\label{sec:org07de76c}
\begin{description}
\item[{Bark, Idbol}] A self-proclaimed "king" threatening travelers and
residents in the \emph{Gloaming}.
\item[{Belcaer}] A trade town in the \emph{Gloaming}.
\item[{Chandlers Guild}] A powerful company with offices across Teraum.

\item[{Crown of Idbol Bark}] A simple "crown" twisted out of discarded
utensils.  \begin{center}\includegraphics[width=.4\linewidth]{../assets/crown-of-idbol-bark.png}\end{center}
\item[{Gloaming}] Sparsely populated, the region is covered in fens broken
by rocky outcroppings.
\item[{goblins}] Goblins are a race of humanoids that vanished 80 years
ago, and certainly don't exist now.
\item[{Guildhouse}] One of the largest buildings in \emph{Belcaer}, home to the
offices of the \emph{Chandlers Guild}, among others.

\item[{Harbook, Abeth}] The primary representative of the \emph{Chandlers
Guild} in the town of \emph{Belcaer}, with an office in
the \emph{Guildhouse}.
\item[{Jacob's Folly}] A limestone tower built long in the past, well outside
the town of \emph{Belcaer}.
\end{description}

\section*{Rules}
\label{sec:org75eb3e7}
\textbf{\emph{Brave Old World}} is a set of rules for telling a \textbf{\emph{story}} by having a
structured conversation. It requires a six-sided dice (two is easier),
a sheet of paper, something to write with, and two to six players
(four to six is better).

One player takes the role of \textbf{\emph{narrator}}, who starts and guides the
story, and the others take the roles of \textbf{\emph{adventurers}}, characters in
that story. The story and characters may be created by the players as
a group, or sourced from pre-written \textbf{\emph{Adventures}}.

The \emph{narrator} says where the adventurers are and what is happening
around them. The adventurers say what it is they'd like to do. When an
adventurer takes a \textbf{\emph{risk}}, they roll two six-sided die (or one dice
twice) and add the results.

\begin{itemize}
\item If they \textbf{roll} a 10 or more (\textbf{10+}), they \textbf{succeed} at what they
attempted.
\item If they \textbf{roll} a 7 through 9 (\textbf{7-9}), they \textbf{partially succeed},
accomplishing what they attempted but with a cost or consequence.
\item If they \textbf{roll} a 6 or less (\textbf{6-}), they \textbf{fail}, and something is
going to happen.
\end{itemize}

If the adventurer is \textbf{\emph{knowledgeable}} about or \textbf{\emph{skilled}} at what
triggered the \emph{roll}, they can add 1 (+1) to their \textbf{\emph{result}}.

After an player \emph{rolls}, the \emph{narrator} says how the game world reacts
to their \emph{result}, and the conversation continues, with players then
saying how their adventurer would respond to those reactions.

This continues until players run out of time or reach a good stopping
point in the \emph{story}.
\end{document}
